% ==============================================================================
% ANÁLISIS DE RESULTADOS (antes "Discusión de resultados")
% ==============================================================================
% Interpreta la evidencia de k-resultados.tex. Solo componente de
% ciberseguridad.
%
% Enfoque (decisión del estudiante, 2026-10-09): la discusión argumenta el
% cumplimiento de los objetivos frente a los criterios de aceptación
% definidos en g-objetivos.tex y h-alcance.tex (Alcance técnico: extracción
% <= 5 ms y Delta p95 <= 15 ms en carga normal y de ataque; validación
% principal en entorno controlado, validación con startups complementaria).
% El escenario de volumen se evalúa contra el mismo criterio de 15 ms y no se cumple. Los hallazgos fuera de esos criterios (falsos
% positivos en aplicaciones de terceros) se presentan como condiciones de
% despliegue y trabajo futuro. Ninguna cifra cambia respecto a Resultados.
%
% "Variantes desconocidas de los vectores cubiertos" (objetivo general): la
% evidencia es la Ronda 5, corpus de payloads de las cuatro categorías
% generados con sqlmap y generadores propios y filtrados por exclusión
% contra el entrenamiento; la detección del sistema se cuenta como bloqueo
% o marca de anomalía (Procedimiento 4.4).
%
% kapoor2023: taxonomía verificada en el PDF del artículo (2026-10-09):
% [L2] "Model uses features that are not legitimate"; [L3] "Test set is not
% drawn from the distribution of scientific interest".
%
% Ronda 4: los 40 ataques que el WAF no bloqueó fueron 11 SQLi, 3 XSS y 26
% Path Traversal, ninguno de Command Injection.
\ifdefined\modoensamble
\else
	\documentclass[11pt, letterpaper]{report}
\usepackage[spanish, es-nodecimaldot, es-noquoting]{babel}
\selectlanguage{spanish}
\usepackage[utf8]{inputenc}
\usepackage[T1]{fontenc}

% ===============================================================================
% El estudiante debe llenar sus datos en esta sección para que la plantilla los 
% auto-importe y genere automáticamente las páginas de portada y de firmas 
% autorizadas.
% ===============================================================================
% Datos del estudiante:
% -------------------------------------------------------------------------------
% Nombre completo
\def \nombreestudiante {Sebastián Huertas Gómez}
% Carné
\def \uvgcarne {22295}
% Facultad
\def \uvgfacultad {Ingeniería}
% Carrera
\def \uvgcarrera {Ingeniería en Ciencia de la Computación y Tecnologías de la Información}

% Datos del trabajo:
% -------------------------------------------------------------------------------
% Título completo (tomado literalmente de protocolo_sebastian_v1 (3).pdf, sección 1)
\def \titulotesis {logSguarDian: Desarrollo de una librería para la detección de amenazas de seguridad en Endpoints de aplicaciones}
% Mes y año de entrega (confirmado por el estudiante: noviembre 2026)
\def \mesentrega {noviembre}
\def \anoentrega {2026}
% Asesor (según protocolo_sebastian_v1 (3).pdf, portada)
\def \nombreasesor {Ing. Erick Francisco Marroquín Rodríguez}
% El catedrático del curso (Ing. Bacilio Alexander Bolaños Lima) figura como
% primer examinador en la hoja de firmas, definido más abajo.

% Datos del tribunal examinador:
% -------------------------------------------------------------------------------
% Nombre del primer examinador (segunda línea de la hoja de firmas)
\def \nombreprimerex {Ing. Ma. Bacilio Alexander Bolaños Lima}

% Capítulos pre-definidos
% -------------------------------------------------------------------------------
% Comentar las líneas de las secciones que desean omitirse, por defecto se
% se incluyen todas.
% Se descomentan conforme se redacta cada capítulo (todavía sin archivo de
% contenido: prefacio, alcance, anexos, glosario).
\def \CAPprefacio {Prefacio}
\def \CAPantecedentes {Antecedentes}
\def \CAPalcance {Alcance}
\def \CAPmetodologia {Metodología}
\def \CAPresultados {Resultados}
\def \CAPdiscusion {Análisis de resultados}
\def \CAPanexos {Anexos}
\def \CAPglosario {Glosario}

% Formato y estilo de la plantilla
% -------------------------------------------------------------------------------
% Modo impresión: Puede des-comentar la siguiente línea para generar un documento pdf sin la portada, para cuando se desee imprimir el documento para encuadernación
% NOTA: la plantilla original traía esta línea activa por defecto, lo cual
% suprime la portada verde (main.tex hace \let\CAPportada\undefined cuando
% \printver está definido). Se comenta para que la versión digital, con
% portada, sea la que se genera por defecto.
%\def \printver {Versión del documento para impresión}

% Portada: Puede cambiarse la imagen en la portada al cambiar el nombre del 
% archivo siguiente. NOTA: debe tener la suficiente resolución para cubrir el área
% designada
\def \imagenportada {plantilla/portadacit.jpg}

% Referencias: Puede des-comentar la siguiente línea para utilizar el formato de referencias APA
%\def \usarAPA {Usar formato APA}

% Párrafo: Puede comentar la siguiente línea si desea emplear un formato de 
% párrafo distinto al establecido por defecto
\def \parpordefecto {Formato de párrafo por defecto}

% Capítulos y secciones: Puede des-comentar la siguiente línea para establecer el
% formato de los capítulos y secciones bajo el estándar original de UVG para
% trabajos de graduación. Este incluye: capítulos con numeración romana, secciones
% con letras mayúsculas, sub-secciones con números y sub-sub-secciones con letras
% minúsculas
%\def \capsecuvg {Formato UVG para capítulos y secciones}
% ================================================================================
% En este archivo se colocan opciones adicionales para modificar el formato de la
% plantilla, para emplearse en otros tipos de documentos que no sean trabajos de
% graduación. Si usted está trabajando su tesis, NO modifique este archivo
% ================================================================================
% Capítulos pre-definidos
% --------------------------------------------------------------------------------
% Comentar las líneas de las secciones que desean omitirse, por defecto se 
% se incluyen todas.
\def \CAPportada {Portada}
\def \CAPcaratula {Caratula}
\def \CAPfirmas {Hoja de firmas}
\def \CAPindice {Índice general}
\def \CAPfiguras {Listado de figuras}
\def \CAPcuadros {Listado de cuadros}
\def \CAPresumen {Resumen}
\def \CAPabstract {Resumen}
\def \CAPintroduccion {Introducción}
\def \CAPobjetivos {Objetivos}
\def \CAPjustificacion {Justificación}
\def \CAPmarcoteorico {Marco teórico}
\def \CAPconclusiones {Conclusiones}
\def \CAPrecomendaciones {Recomendaciones}
\def \CAPbibliografia {Bibliografía}

\usepackage{xcolor}
\usepackage{amsfonts}
\usepackage{amsmath}
\usepackage{amssymb}
\usepackage{amsthm}
\usepackage{mathtools}
\usepackage{graphicx}
\usepackage{xfrac}
\usepackage{float}
\usepackage[hypertexnames=false]{hyperref}
\usepackage[font=small]{caption}
\usepackage{subcaption}
% Cuadros: array (columnas L{} alineadas a la izquierda) y longtable (cuadros
% que ocupan más de una página). Se cargan aquí, no en
% 2-paquetes_y_comandos_usuario.tex, para que los capítulos de Metodología y
% Resultados no dependan de ese archivo.
\usepackage{array}
\usepackage{longtable}
\newcolumntype{L}[1]{>{\raggedright\arraybackslash}p{#1}}
\renewcommand{\arraystretch}{1.3}
\setlength{\LTpre}{0.3cm}
\setlength{\LTpost}{0.1cm}
% Margen de estiramiento para que los párrafos con palabras largas no
% generen líneas que sobresalgan del margen.
\setlength{\emergencystretch}{3em}
\usepackage{xpatch}
\usepackage{emptypage}
% htt: permite dividir con guion el texto en \texttt (nombres de archivo,
% clases y campos), que de otro modo sobresale del margen derecho.
\usepackage[htt]{hyphenat}
\usepackage{fancyhdr}
\usepackage[backend=biber, style=ieee]{biblatex}
\ifdefined\usarAPA
    \usepackage[backend=biber, style=apa]{biblatex}
\fi
\addbibresource{m-bibliografia.bib}

\usepackage[percent]{overpic}
\usepackage{chngcntr}

\ifdefined\CAPglosario
	\usepackage[numberedsection]{glossaries}
	\makeglossaries
    \newglossaryentry{rasp}
{
    name={RASP},
    description={Autoprotección de aplicaciones en tiempo de ejecución (\textit{Runtime Application Self-Protection}): mecanismo de seguridad que se ejecuta dentro de la propia aplicación y analiza sus peticiones para detectar y bloquear ataques.}
}

\newglossaryentry{waf}
{
    name={WAF},
    description={Cortafuegos de aplicaciones web (\textit{Web Application Firewall}): dispositivo o servicio perimetral que filtra el tráfico HTTP antes de que llegue a la aplicación.}
}

\newglossaryentry{modsecurity}
{
    name={ModSecurity},
    description={Motor de WAF de código abierto que se usa con reglas externas; en este trabajo se configuró con OWASP CRS.}
}

\newglossaryentry{crs}
{
    name={OWASP CRS},
    description={Conjunto de reglas genéricas de OWASP (\textit{Core Rule Set}) para ModSecurity, con niveles de paranoia que controlan qué tan estrictas son.}
}

\newglossaryentry{owasp}
{
    name={OWASP Top 10},
    description={Lista de las diez categorías de riesgo más críticas para aplicaciones web, publicada por OWASP; este trabajo usa la edición 2021.}
}

\newglossaryentry{mitre}
{
    name={MITRE ATT\&CK},
    sort={MITRE ATT&CK},
    description={Base de conocimiento de tácticas y técnicas de atacantes, usada aquí para asociar cada vector de ataque a una técnica.}
}

\newglossaryentry{cwe}
{
    name={CWE},
    description={Enumeración común de debilidades de software (\textit{Common Weakness Enumeration}): catálogo de tipos de defectos que originan vulnerabilidades.}
}

\newglossaryentry{sqli}
{
    name={SQL Injection},
    description={Ataque que inserta sentencias SQL en una entrada de la aplicación para alterar las consultas que ejecuta la base de datos.}
}

\newglossaryentry{xss}
{
    name={XSS},
    description={\textit{Cross-Site Scripting}: ataque que inyecta código ejecutable en una página que otros usuarios verán.}
}

\newglossaryentry{pathtraversal}
{
    name={Path Traversal},
    description={Ataque que manipula rutas de archivo (por ejemplo con \texttt{../}) para leer archivos fuera del directorio previsto; incluye la inclusión local de archivos (LFI).}
}

\newglossaryentry{cmdi}
{
    name={Command Injection},
    description={Ataque que logra ejecutar comandos del sistema operativo a través de una entrada de la aplicación.}
}

\newglossaryentry{middleware}
{
    name={\textit{middleware}},
    sort={middleware},
    description={Función que se ejecuta entre la recepción de una petición y la lógica de negocio de la aplicación; en Express se registra con \texttt{app.use}.}
}

\newglossaryentry{rf}
{
    name={Random Forest},
    description={Algoritmo supervisado que combina muchos árboles de decisión; en la librería es el único modelo con autoridad para bloquear una petición.}
}

\newglossaryentry{iforest}
{
    name={Isolation Forest},
    description={Algoritmo no supervisado de detección de anomalías que aísla las observaciones atípicas con particiones aleatorias; en la librería solo registra anomalías y nunca bloquea.}
}

\newglossaryentry{onnx}
{
    name={ONNX},
    description={Formato abierto para serializar modelos de aprendizaje automático, que permite ejecutarlos fuera de Python.}
}

\newglossaryentry{workerthreads}
{
    name={\texttt{worker\_threads}},
    sort={worker_threads},
    description={Módulo de Node.js que ejecuta código en hilos independientes del hilo principal; la inferencia de los modelos se ejecuta ahí.}
}

\newglossaryentry{eventloop}
{
    name={\textit{Event Loop}},
    sort={Event Loop},
    description={Bucle de eventos de Node.js que atiende las peticiones en un único hilo; bloquearlo degrada a toda la aplicación.}
}

\newglossaryentry{failopen}
{
    name={\textit{fail-open}},
    sort={fail-open},
    description={Política por la cual, ante un error o un tiempo de espera del módulo de detección, la petición se deja pasar en lugar de bloquearse.}
}

\newglossaryentry{monitor}
{
    name={Modo \texttt{monitor}},
    sort={Modo monitor},
    description={Modo de operación en el que la librería registra las peticiones que bloquearía, pero no las bloquea.}
}

\newglossaryentry{block}
{
    name={Modo \texttt{block}},
    sort={Modo block},
    description={Modo de operación en el que la librería responde HTTP 403 a las peticiones que clasifica como ataque.}
}

\newglossaryentry{pesp}
{
    name={p95},
    description={Percentil 95 del tiempo de respuesta: el valor por debajo del cual queda el 95~\% de las peticiones.}
}

\newglossaryentry{deltap95}
{
    name={$\Delta p95$},
    sort={Delta p95},
    description={Diferencia entre el p95 con el middleware activo y el p95 de la línea base sin él; en este trabajo se mide en milisegundos.}
}

\newglossaryentry{funo}
{
    name={F1},
    description={Media armónica de la precisión y el recall; es la métrica principal de clasificación del trabajo.}
}

\newglossaryentry{auc}
{
    name={AUC-ROC},
    description={Área bajo la curva ROC; mide qué tan bien separa un modelo dos clases sin depender de un umbral.}
}

\newglossaryentry{wilson}
{
    name={Intervalo de Wilson},
    description={Intervalo de confianza para una proporción, adecuado para muestras pequeñas o proporciones cercanas a 0 o 1.}
}

\newglossaryentry{casidup}
{
    name={Casi-duplicado},
    description={Registro muy parecido a otro pero no idéntico; si queda repartido entre entrenamiento y prueba, puede inflar las métricas.}
}

\newglossaryentry{minhash}
{
    name={MinHash},
    description={Técnica para estimar rápidamente la similitud entre conjuntos, usada para detectar casi-duplicados.}
}

\newglossaryentry{fuga}
{
    name={Fuga de datos},
    description={Presencia en el conjunto de evaluación de información disponible durante el entrenamiento, lo que sobreestima el desempeño.}
}

\newglossaryentry{corpuslimpio}
{
    name={Corpus limpio},
    description={Conjunto de payloads de evaluación filtrado para excluir los que coinciden con el entrenamiento; en este trabajo es el de la Ronda~5.}
}

\newglossaryentry{useragent}
{
    name={\texttt{User-Agent}},
    sort={User-Agent},
    description={Cabecera HTTP con la que el cliente se identifica (navegador, cURL, entre otros); su representación en el corpus sesgó al modelo.}
}

\newglossaryentry{falsopositivo}
{
    name={Falso positivo},
    description={Petición legítima que el sistema clasifica o bloquea como ataque.}
}

\newglossaryentry{parmodelos}
{
    name={Par de modelos},
    description={Combinación de un archivo \texttt{rf.onnx} y uno \texttt{if.onnx} identificada por su suma SHA-256 (pares A, B y C).}
}

\newglossaryentry{tarball}
{
    name={\textit{tarball}},
    sort={tarball},
    description={Archivo comprimido \texttt{.tgz} de un paquete de npm, publicado o construido localmente.}
}

\newglossaryentry{kanban}
{
    name={Kanban},
    description={Metodología de gestión basada en un tablero de columnas que representan el estado de las tareas.}
}

\newglossaryentry{artillery}
{
    name={Artillery},
    description={Herramienta de pruebas de carga usada para medir la latencia con y sin el middleware.}
}

\newglossaryentry{sqlmap}
{
    name={sqlmap},
    description={Herramienta de código abierto que automatiza la detección y la explotación de inyecciones SQL; generó los payloads de SQL Injection del corpus limpio.}
}

\newglossaryentry{zap}
{
    name={OWASP ZAP},
    description={Herramienta de OWASP para pruebas de seguridad de aplicaciones web; se usó para confirmar el payload de XSS.}
}

\newglossaryentry{mcp}
{
    name={MCP},
    description={Protocolo de Contexto de Modelo (\textit{Model Context Protocol}): estándar para que un asistente de inteligencia artificial use herramientas externas.}
}

\fi

\usepackage[top=1in, left=1.5in, right=1in, bottom=1in]{geometry}
\usepackage[Sonny]{fncychap}

\graphicspath{ {figuras/} }
\definecolor{uvg-green}{RGB}{17,71,52}
\newcommand{\defaultparformat}[1]{
	{\setlength{\parskip}{2ex}
     \input{#1}}
}
\ifdefined\capsecuvg
	\renewcommand\thechapter{\Roman{chapter}}
    \renewcommand\thesection{\Alph{section}}
	\renewcommand\thesubsection{\arabic{subsection}}
    \renewcommand\thesubsubsection{\alph{subsubection}}
\fi
\counterwithout{figure}{chapter}
\counterwithout{table}{chapter}
\counterwithout{equation}{chapter}

\providecommand{\blankpage}{
\newpage
\thispagestyle{empty}
\mbox{}
\newpage
}

% ==============================================================================
% Paquetes y comandos adicionales del usuario
% ==============================================================================
% TikZ: requerido por las figuras del capítulo de Marco Teórico (pipeline de
% middleware, comparación RASP vs. control perimetral).
\usepackage{tikz}
\usetikzlibrary{arrows.meta, positioning, shapes.geometric, fit, backgrounds}


\pagestyle{headings}
\begin{document}

\ifdefined\parpordefecto
	\setlength{\parskip}{2ex}
\fi

\fi

\chapter{Análisis de resultados}
\label{cap:discusion}

Este capítulo interpreta la evidencia presentada en el Capítulo~\ref{cap:resultados} frente a los
criterios de aceptación definidos en los Objetivos y en el Alcance. Se analiza, objetivo por
objetivo, qué significan los resultados y en qué medida sustentan su cumplimiento; se contrastan
con los trabajos previos y los fundamentos del Marco Teórico; se analizan las condiciones de
despliegue y las limitaciones que acotan los hallazgos; se valida la solución frente a sus
criterios de aceptación; y se cierra con sugerencias que se desprenden de los resultados y que
dan paso a las conclusiones y a las recomendaciones. En cada apartado se distingue lo comprobado
de lo que es solo una explicación posible.

% ------------------------------------------------------------------------------
\section{Objetivo general: una librería RASP funcional para Node.js}

El objetivo general establece cuatro componentes verificables: una librería npm, un pipeline de
ingeniería de características, un modelo híbrido capaz de detectar y clasificar amenazas conocidas
y variantes desconocidas de los vectores cubiertos, y la operación en tiempo de ejecución sin
comprometer el rendimiento del servidor. Los resultados sustentan los tres primeros en lo esencial
y solo en parte el cuarto, como se detalla al final de esta sección.

\textbf{Librería y pipeline.} \texttt{logsguardian} está publicada en el registro público de
npm (0.1.0 y 0.1.1), se instala con un solo comando y no requiere infraestructura externa. Integra un extractor de
76 características implementado una sola vez en TypeScript, de modo que el entrenamiento y la
producción calculan el mismo vector, y un modelo híbrido ejecutado en hilos de trabajo que no
bloquean el Event Loop. La paridad de $\leq2.42\times10^{-7}$ entre los modelos de Python y sus
versiones ONNX garantiza que la serialización no altera las predicciones. Las evaluaciones
emplearon tres pares de modelos con la misma arquitectura de ejecución
(Cuadro~\ref{tab:pares-modelos}): el paquete 0.1.0 contiene \texttt{rf\_v10}/\texttt{if\_v9}, que
consume 67 y 61 características, y las métricas del conjunto de prueba corresponden a
\texttt{rf\_v11}/\texttt{if\_v10}, que consume 69 y 63 y es el contenido del paquete 0.1.1,
publicado el 11 de octubre de 2026 con la corrección del middleware. Como se discute en la
Sección~\ref{sec:dis-oe3}, las generaciones evaluadas cumplen el criterio de cobertura sobre
payloads no vistos, aunque con el paquete 0.1.0 el margen en Path Traversal es estrecho. Con la 0.1.1 se midió la cobertura de las Configuraciones 2 y 3b de la Ronda~5; las aplicaciones de
terceros y el rendimiento no se midieron con ella, de modo que lo que el paquete vigente haría en
esas pruebas queda por medir. Este resultado cubre la brecha identificada en los Antecedentes: las soluciones
RASP comerciales requieren licencias o plataformas externas, y la única alternativa abierta para
Node.js (node-rasp) exigía modificar el runtime y fue abandonada \cite{saparchive2019}.

\textbf{Amenazas conocidas.} La detección y clasificación de las cuatro categorías se sustenta en
el OE3 (Sección~\ref{sec:dis-oe3}): F1 de entre 0.9239 y 0.9956 sobre el conjunto de prueba.

\textbf{Variantes desconocidas de los vectores cubiertos.} La evidencia de esta capacidad es el
corpus limpio de la Ronda~5: 600 payloads de SQL Injection, Path Traversal y Command Injection
generados con sqlmap y con generadores propios, y filtrados para excluir toda coincidencia exacta
del campo completo con el entrenamiento. Son, por construcción, variantes que no coinciden de forma
exacta con las que el modelo vio (4 de los 200 payloads de SQL Injection coinciden con alguna
partición por igualdad normalizada, según la medición de solapamiento del trabajo de graduación de Diego Valenzuela \cite{valenzuela2026ia}, y ninguno de Path Traversal ni de Command Injection), con técnicas de codificación (doble codificación URL, UTF-8 sobrelargo, barra invertida,
byte nulo) y separadores de shell distintos de los del entrenamiento. El paquete 0.1.0 bloqueó
el 100~\% de las de SQL Injection, el 91.5~\% de las de Command Injection y el 85.0~\% de las de
Path Traversal, y el modelo vigente, el 99.5~\%, el 96.0~\% y el 95.0~\%. Esta capacidad corresponde
al Random Forest: la cobertura cuenta solo peticiones bloqueadas y el Isolation Forest no bloquea.
El Isolation Forest cumple sus dos criterios sobre el conjunto de prueba (91.25~\% de detección con
5.46~\% de falsos positivos), pero su contribución a la detección de variantes no vistas solo tiene
evidencia indirecta: en el conjunto de validación del componente de Inteligencia Artificial \cite{valenzuela2026ia} marca 19
de los 39 ataques que el Random Forest deja pasar, y bloquear sus marcas agregaría 866 benignos
bloqueados; esas cifras son optimistas, porque el umbral del Isolation Forest se eligió sobre ese
mismo conjunto, y no se evaluaron sobre el conjunto de prueba. Con el paquete 0.1.0, el límite inferior del intervalo de confianza de Path Traversal
(79.4~\%) quedaba por debajo del criterio, pero esa cifra incluía dos defectos de medición ya
corregidos; con el paquete 0.1.1 y el ejecutor corregido es de 100~\% (98.1--100.0~\%), medido
sobre el mismo corpus que reveló los defectos y, por tanto, sin independencia
(Sección~\ref{sec:met-validez}). El corpus propio de 200 payloads de XSS, medido con el paquete 0.1.1, se bloqueó en su
totalidad (200 de 200; intervalo de Wilson al 95~\%: 98.1--100.0~\%); se construyó después de
conocer los resultados de la Ronda~5, y 3 de sus payloads coinciden con el entrenamiento por
subcadena normalizada, de modo que no es una evaluación ciega. Con esas salvedades, el sistema
reconoce variantes no vistas de los vectores cubiertos en una proporción superior al criterio de
cobertura del 80~\%.

\textbf{Rendimiento del servidor anfitrión.} El Alcance fija dos criterios cuantitativos: la
extracción de características en 5~ms o menos por petición y un incremento de la latencia p95 de
15~ms como máximo bajo carga normal y bajo carga de ataque. Ambos se cumplen: la extracción,
medida con la ida y vuelta completa a su \texttt{worker\_thread}, tardó 0.0273~ms en el p95, más
de 180 veces por debajo del límite, y el incremento de latencia fue de
12.08~ms y 12.85~ms. El valor relativo de esos mismos incrementos (+485~\% y +278~\%) confirma por
qué el criterio se expresó en milisegundos: la aplicación de referencia responde en 2 a 5~ms, de
modo que cualquier costo fijo produce porcentajes elevados sin que el usuario perciba una
degradación proporcional; en una API con tiempos de respuesta de decenas o cientos de
milisegundos, 12~ms adicionales representan una fracción menor. Con ello, el criterio de latencia se cumple en carga normal y de ataque; en el escenario de
volumen no se cumple (\(\Delta p95\) de 6{,}495.46~ms y 7.00~\% de tiempos de espera).

\textbf{Costo de recursos.} El protocolo definió además dos indicadores de costo, que no se
cumplen: un incremento de memoria de 150~MB (143.1~MiB) como máximo y un incremento de CPU de
15~\% como máximo. Con el middleware activo, el contenedor usó 332.0~MiB más de memoria bajo carga
normal y 331.9~MiB más bajo carga de ataque, más del doble del límite, y el uso de CPU aumentó en
89 y 85 puntos porcentuales. El incremento de memoria es casi idéntico entre escenarios, por lo
que corresponde a un costo fijo de la arquitectura al activarse (el runtime nativo de ONNX y tres
hilos de inferencia con sus modelos cargados) y no a un crecimiento con la carga. El límite de
150~MB se fijó antes de medir la memoria nativa que reserva \texttt{onnxruntime-node} al cargar
modelos de árboles, proporcional a su número de nodos, y pensando en instancias de 1~GB; en una
instancia de ese tamaño, 332~MiB representan cerca de un tercio de la memoria disponible, un costo
que debe considerarse al dimensionar la infraestructura. El incremento de CPU, medido en un
ejecutor de 2~vCPU a 100 peticiones por segundo, equivale aproximadamente a un núcleo dedicado a
la inferencia. Ambos costos son objetivos directos de optimización para la siguiente versión.

\textbf{Veredicto sobre el objetivo general.} El objetivo general se cumple de forma parcial. Se
cumplen la librería publicada, el pipeline, la clasificación de las cuatro categorías, la cobertura
de variantes no vistas con el Random Forest (con las reservas indicadas) y el criterio de latencia
en carga normal y de ataque. No se cumplen los indicadores de memoria y CPU, ni el criterio de
latencia en el escenario de volumen, y el aporte del Isolation Forest a la detección de variantes
desconocidas solo tiene evidencia indirecta. Atribuir ese resultado a la arquitectura y no a un
modelo en particular es razonable para el costo de recursos, que es casi constante entre
escenarios, pero no está comprobado: el costo por componente no se midió.

% ------------------------------------------------------------------------------
\section{Objetivo específico 1: selección de vectores}

Los cuatro vectores quedaron mapeados a OWASP Top 10 2021 y a MITRE ATT\&CK con fuentes
oficiales verificadas, y ese mapeo cumplió la función que el objetivo le asigna: definió las cinco
clases canónicas del dataset y las categorías sobre las que se calculan todas las métricas del
OE3. La precisión del mapeo se reforzó durante el proyecto: Path Traversal/LFI se ubicó en
A01:2021 (Broken Access Control) y no en A03:2021 (Injection), lo que muestra que la selección de
los cuatro vectores responde a un criterio operativo (todos materializan su payload en la petición
HTTP y son detectables con las mismas características) y no a su pertenencia a una sola categoría
de OWASP. El OE1 se cumple.

% ------------------------------------------------------------------------------
\section{Objetivo específico 2: dataset etiquetado}

El corpus de 383{,}202 peticiones etiquetadas supera 3.8 veces el mínimo de 100{,}000 muestras. Su
representatividad se sustenta en la diversidad de origen: siete fuentes públicas auditadas (más de
la mitad de las filas proviene de la colección CAPEC y una cuarta parte de modsec-learn), entre
ellas registros reales de un servidor en producción protegido por ModSecurity \cite{lucz2025}
(1.0~\% de las filas), complementadas con dos listas públicas de payloads y con conjuntos propios,
con las cinco clases presentes en las tres particiones. Construir el corpus a partir de fuentes
públicas, en lugar de generarlo en su totalidad en el entorno controlado, aportó un volumen y una
variedad de técnicas de ataque que el entorno no habría podido producir en el tiempo disponible;
el entorno controlado se empleó para la validación, que es donde su reproducibilidad aporta más
valor. El desbalance de clases se trató con ponderación en el entrenamiento, y su efecto queda
acotado por los resultados: incluso la clase minoritaria, Command Injection (2.6~\% del conjunto de
prueba), alcanzó un F1 de 0.9239. El objetivo pide balancear tráfico legítimo y ataques: el corpus contiene
ambos en las cinco clases y en las tres particiones, pero no los equilibra en conteos (el 59.3~\%
del conjunto de prueba es SQLi), de modo que el criterio se cumple como presencia de ambos tipos de
tráfico y no como proporciones iguales. El etiquetado determinista y la partición verificada por hash
permiten reconstruir cada partición a partir de sus fuentes públicas, con la excepción de los
conjuntos sintéticos de Command Injection, cuyo generador no está versionado. El OE2 se cumple.

% ------------------------------------------------------------------------------
\section{Objetivo específico 3: efectividad de la detección}
\label{sec:dis-oe3}

El criterio del OE3 exige un F1 de al menos 0.80 en tres de las cuatro categorías. El sistema lo
alcanzó en las cuatro, con un margen mínimo de 0.12 sobre el criterio (Command Injection, 0.9239) y
un F1 macro de 0.9776. El AUC-ROC de entre 0.99899 y 0.99982 en las categorías de ataque indica que la capacidad de discriminar
cada categoría no depende de un umbral particular. La precisión de 0.9984 en la clase legítima
muestra que, dentro de la distribución del corpus, el modelo cumple el balance entre falsos
negativos y falsos positivos que el objetivo plantea para un entorno sin analista dedicado.

La cobertura sobre payloads no vistos confirma que estas métricas no dependen de memorizar el
entrenamiento. La auditoría de fuga fue decisiva para esa conclusión: en Command Injection, el
100~\% de la Ronda~4 correspondía en parte a plantillas presentes en el entrenamiento, y la cifra
atribuible a detección real, medida con el corpus limpio, es 91.5~\% con el paquete 0.1.0 y
96.0~\% con el tarball del modelo vigente, y 92.5~\% con el paquete 0.1.1 y el ejecutor corregido.
Esta última cifra es más baja que la anterior, pero no significa que la detección haya empeorado:
el ejecutor anterior volvía a codificar los payloads con separadores ya codificados, y 10 de ellos,
corrompidos así, se bloqueaban, mientras que enviados correctamente no se bloquean (HTTP 200); es la
primera medición correcta de esa técnica de evasión, y un límite real de la librería. En Path
Traversal, el paquete 0.1.0 alcanzó 85.0~\% con un intervalo de confianza cuyo límite inferior
(79.4~\%) quedaba por debajo del criterio, pero esa cifra no describe solo al paquete: sus fallos se
concentraban en las categorías de codificación y tenían dos causas independientes, un defecto del
middleware, que añadía una capa de codificación a la consulta, y un defecto del ejecutor de pruebas,
que volvía a codificar payloads ya codificados. Ninguna corrección bastó por sí sola: con el
middleware corregido y el ejecutor anterior los fallos fueron los mismos 10 que sin la corrección, y
desaparecieron al corregir también el ejecutor. Con el paquete 0.1.1 instalado desde el registro de
npm, la cobertura de Path Traversal es de 200 de 200 (intervalo de Wilson al 95~\%: 98.1--100.0~\%),
por encima del criterio con un margen amplio. Esa cifra tiene dos límites. No se midió el paquete
0.1.0 con el ejecutor corregido, por lo que no puede repartirse la mejora entre el modelo, el
middleware y el ejecutor. Y se obtuvo sobre el mismo corpus que sirvió para encontrar los
defectos, de modo que es optimista por construcción y debe confirmarse con un corpus nuevo. En XSS,
la detección se sustenta en el F1 de 0.9877 sobre el conjunto de prueba, en el 98~\% de la suite
extremo a extremo, en el 97.3~\% de la Ronda~4 y, desde el 25 de octubre de 2026, en la cobertura
del corpus propio de 200 payloads: 200 de 200 (intervalo de Wilson al 95~\%: 98.1--100.0~\%) y 64
de 64 en el campo vulnerable \texttt{bio} (94.3--100.0~\%). El criterio de cobertura del 80~\% se
cumple para XSS con evidencia estadística. Esa cobertura tiene límites declarados: el corpus se
construyó después de conocer los resultados de la Ronda~5, 136 de sus 200 payloads apuntan a
campos sin vulnerabilidad confirmada y 3 coinciden con el entrenamiento por subcadena normalizada.
La explotabilidad confirmada sigue descansando en el único caso que ZAP confirmó sobre la
aplicación de referencia (\texttt{POST /profile}, bloqueado por todos los paquetes evaluados), con
n=1. El OE3 se cumple, sin reservas en las métricas de clasificación y con reservas en la
cobertura: que la medición vigente de Path Traversal y de Command Injection no es independiente
del corpus que motivó las correcciones, y que la de XSS se hizo sobre un corpus construido después
de conocer los resultados de la Ronda~5.

% ------------------------------------------------------------------------------
\section{Defensa en profundidad}
\label{sec:dis-defensa}

En la Ronda~4, ModSecurity con OWASP CRS en nivel de paranoia 1 bloqueó el 93.2~\% de los ataques,
\textit{logSguarDian} el 98.8~\% y ambas capas el 99.5~\%; de los 40 ataques que el WAF dejó pasar,
la librería bloqueó 37. Ese corpus tiene fuga confirmada en Command Injection, pero ninguno de los
40 ataques que el WAF no bloqueó pertenece a esa categoría (fueron 11 de SQL Injection, 3 de XSS y
26 de Path Traversal). Por igualdad normalizada, 7 de los 113 payloads de XSS y ninguno de los 77
de SQL Injection ni de los 200 de Path Traversal coinciden con alguna partición; por subcadena,
26 de SQL Injection y 13 de XSS, un criterio permisivo que el trabajo de graduación de Diego
Valenzuela \cite{valenzuela2026ia} (sección VIII.3.5) califica de límite superior. Con esa salvedad, el resultado apunta al
posicionamiento de \textit{logSguarDian} definido en el Marco Teórico: una capa complementaria al
WAF, que detiene la mayor parte de lo que supera la primera capa, y no un sustituto.

La comparación sobre el corpus limpio de la Ronda~5, que no tiene la fuga de la Ronda~4, se repitió
el 25 de octubre de 2026 con las cuatro configuraciones sobre el mismo árbol de la aplicación, con
el ejecutor corregido y con 800 payloads, 200 de ellos de XSS. El WAF solo y la librería sola tienen
perfiles distintos. En SQL Injection son equivalentes (99.5~\% ambos, 199 de 200). En Path
Traversal no hay ventaja del WAF: con el paquete 0.1.1 la librería bloqueó 200 de 200 y el WAF
197 de 200, una diferencia de tres payloads con intervalos de Wilson solapados. En Command
Injection la librería bloquea más (92.5~\% frente a 82.0~\%); la cifra del WAF, con intervalo de
Wilson al 95~\% de 76.1 a 86.7~\%, no permite afirmar que supere el 80~\%. El 68.0~\% que se había
reportado para el WAF era un efecto del ejecutor anterior: al volver a ejecutarlo sobre el mismo
árbol reprodujo 136 de 200, y el ejecutor corregido dio 164. Aun así, el WAF bloquea menos los
payloads de Command Injection con separadores ya codificados (76.2~\%; intervalo de Wilson al
95~\%: 67.2--83.3~\%) que los no codificados (88.4~\%; 80.4--93.4~\%), una diferencia de 12.2
puntos porcentuales que se mantiene con el ejecutor corregido (prueba exacta de Fisher bilateral,
$p = 0.028$; intervalo de confianza de Wald al 95~\% de la diferencia, de 1.8 a 22.6 puntos). Los
dos subconjuntos no están controlados por otras características del payload, por lo que la
diferencia no se atribuye a la codificación como causa. En XSS, el WAF bloqueó 189 de 200 (94.5~\%) y falló en 11 payloads con doble
codificación URL enviados como JSON o \texttt{multipart/form-data}; la librería bloqueó los 200.

En conjunto, la librería sola bloqueó 784 de 800 payloads, el WAF solo 749 y la combinación 797: la
combinación agrega 13 payloads a la librería (1 de SQL Injection y 12 de Command Injection) y 48
al WAF. De los 51 ataques que el WAF no bloqueó, la combinación bloqueó 48 y la librería sola 44. La
complementariedad es asimétrica: en Path Traversal y en XSS la librería sola bloquea todo lo que
bloquea la combinación, y solo en SQL Injection y en Command Injection cada capa detiene payloads
que la otra deja pasar. Por eso no puede afirmarse que la combinación supere a ambas capas en todas
las categorías, ni que ninguna subsuma a la otra. Cuatro payloads de Command Injection los bloquea
la combinación y no los bloquean ni la librería sola ni el WAF solo; las cuatro cabeceras
\texttt{X-*} que agrega el proxy lo explican: al enviar esos payloads a la Configuración~2 con las
cabeceras, la librería los bloqueó en las tres repeticiones, y sin ellas ninguno, porque las
cabeceras elevan de 0 a 4 el recuento de cabeceras fuera de la lista estándar que usa el Random
Forest. La Configuración~2 y la librería dentro de la 3b no son, por tanto, la misma medición, y el
veredicto de la librería depende en parte de las cabeceras que agrega el despliegue. Tres payloads
de Command Injection (\texttt{cmdi\_6}, \texttt{cmdi\_184} y \texttt{cmdi\_185}) no los bloquea
ninguna configuración: el WAF no activa ninguna regla de CRS sobre ellos (puntaje de entrada 0
frente al umbral de 5) y, en la Configuración~2, la clase de mayor probabilidad del Random Forest
es benigna en \texttt{cmdi\_6} y, con las cabeceras del proxy, de ataque con probabilidad 0.3333 en
\texttt{cmdi\_184} y \texttt{cmdi\_185}, a 0.0167 del umbral de 0.35 (a 0.05 en \texttt{cmdi\_185}
sin cabeceras); no se investigó por qué. Las Rondas 1 a 3 aportaron
además un requisito de integración: la librería debe montarse después de los analizadores de
cuerpo de la aplicación, porque con \texttt{multipart/form-data} procesado por ruta no veía el
contenido de la petición.

% ------------------------------------------------------------------------------
\section{Contraste con los antecedentes y el marco teórico}

Las métricas de clasificación de este trabajo están al nivel de las reportadas en los
Antecedentes: F1 de 0.9956 en SQL Injection frente a un F1 de 0.98 del Random Forest de Peralta-Garcia et al.
\cite{peralta2024sql}, y F1 de 0.9239 en Command Injection frente a una exactitud de 99.3~\% y un
\textit{recall} de 98.2~\% de Wang et al. \cite{seyyar2024}\footnote{Wang et al. reportan 99.3~\% en
el resumen y 99.21~\% en la comparación de la sección de resultados.}, cuyo modelo de aprendizaje
profundo se evaluó en una GPU y no reporta latencia de inferencia ni despliegue como middleware. La comparación numérica directa tiene límites, porque los
conjuntos de datos, las métricas y las unidades de análisis difieren. Las contribuciones de este
trabajo están en lo que los anteriores no reportan. Frente al WAF de Durmuşkaya y Bayraklı \cite{durmuskaya2025waf}, el más cercano en arquitectura, \textit{logSguarDian} opera dentro del proceso sin un salto de red
adicional; y, a diferencia de los trabajos revisados en los Antecedentes, este proyecto mide la latencia agregada
bajo carga sostenida, evalúa la detección sobre payloads filtrados contra el entrenamiento y prueba
la librería publicada en aplicaciones de terceros.

Frente al Marco Teórico, los resultados confirman dos de sus planteamientos y obligan a matizar un
tercero. Confirman el posicionamiento de un RASP como capa complementaria al WAF y no como su
sustituto: la librería, que analiza la petición dentro del proceso, y el WAF perimetral detectan
conjuntos parcialmente distintos (Sección~\ref{sec:dis-defensa}). Confirman también la descripción
del Isolation Forest como detector que modela el tráfico legítimo y marca lo atípico: detecta el
91.25~\% de los ataques del conjunto de prueba con 5.46~\% de falsos positivos. Obligan a matizar
la expectativa de que ese detector aporte por sí solo a las variantes desconocidas, porque sus
marcas sobre ataques que el Random Forest deja pasar existen (19 de 39 en validación \cite{valenzuela2026ia}) pero
costarían 866 bloqueos indebidos, y por eso solo registra. La advertencia de Kapoor y Narayanan
\cite{kapoor2023} sobre características no legítimas y conjuntos de prueba de la misma distribución
que el entrenamiento resultó pertinente: la cabecera \texttt{User-Agent} y la saturación de
casi-duplicados son ejemplos concretos en este corpus (Secciones~\ref{sec:dis-despliegue} y
\ref{sec:dis-limitaciones}).

% ------------------------------------------------------------------------------
\section{Condiciones de despliegue}
\label{sec:dis-despliegue}

Dos evaluaciones adicionales, una de ellas con un criterio que no se cumple, delimitan las
condiciones en que debe desplegarse.

\textbf{Volumen de datos.} En un escenario no contemplado en el protocolo, pero evaluado contra el mismo criterio de 15~ms, con la base de datos
sembrada y una ruta de listado que ya tardaba 63.84~ms sin middleware, el 7~\% de las peticiones
terminó por tiempo de espera. Los tiempos de espera se concentraron al final de cada repetición,
un patrón consistente con una cola que crece cuando el tiempo de servicio supera la tasa de
llegada. El criterio de 15~ms no se cumple en este escenario ($\Delta p95$ de 6{,}495.46~ms), lo que
limita el cumplimiento del objetivo general a las condiciones de carga normal y de ataque con
respuestas de la aplicación cortas; identifica el límite operativo de la versión actual y
orienta su siguiente iteración.

\textbf{Calibración con tráfico real.} En tres aplicaciones de terceros, con la configuración por
defecto, la librería bloqueó entre el 92.5~\% y el 100~\% del tráfico legítimo probado con un
\texttt{User-Agent} de navegador, con muestras pequeñas o poco diversas (5 peticiones en DVNA; 50
dirigidas a una sola ruta en node-api-goat); en Juice Shop, en modo \texttt{monitor}, el modelo asignó
veredicto de bloqueo al 97.1~\%, y en el estudio con navegadores reales se bloquearon 13 de 15
combinaciones de \texttt{User-Agent} (Cuadro~\ref{tab:res-terceros-ua}). La exclusión por fuente del trabajo del componente de Inteligencia Artificial
\cite{valenzuela2026ia}, con modelos distintos del evaluado, apunta en el mismo
sentido: sin \texttt{modsec\_learn}, \texttt{payloads\_csv} o \texttt{xss\_dataset} en el entrenamiento, el
bloqueo de sus filas legítimas sube de 0.1~\%, 0.4~\% y 0.2~\% a 76.0~\%, 27.0~\% y 94.5~\%. El Alcance definió la validación en el entorno
controlado como criterio principal y la validación externa como escenario complementario; estas
pruebas, realizadas por recomendación del asesor, cumplen ese papel complementario y aportan el
dato que el entorno controlado no podía dar. La hipótesis principal es un sesgo de representación
del corpus: solo el 0.4~\% de las filas legítimas tiene una cabecera \texttt{User-Agent} de más de
40 caracteres, frente al 33.3--97.7~\% de las clases de ataque, y esa característica es la de mayor
contribución en el modelo vigente (\texttt{rf\_v11}) según el análisis SHAP \cite{lundberg2017unified,valenzuela2026ia}, y lo era también en
\texttt{rf\_v3}. Esa hipótesis no es la única explicación: con el \texttt{User-Agent} por defecto de
Node.js, dvws-node bloqueó el 75~\% del tráfico legítimo, y Juice Shop el 39.4~\% con cURL, de modo
que el \texttt{User-Agent} amplifica el bloqueo sin explicarlo por completo, y no se verificó con un
corpus legítimo alternativo. Kapoor y Narayanan describen este tipo
de característica, que codifica el origen de los datos y no la propiedad de interés, como no
legítima, y advierten que un conjunto de prueba de la misma distribución que el entrenamiento no
representa la distribución de interés \cite{kapoor2023}. Si la hipótesis se confirma, el hallazgo sería corregible con datos y
calibración, sin cambiar la arquitectura, y define la configuración recomendada para la población
objetivo: instalación en modo \texttt{monitor} y calibración del umbral con el tráfico propio antes
de activar el bloqueo.

% ------------------------------------------------------------------------------
\section{Cumplimiento de los objetivos}

El Cuadro~\ref{tab:dis-cumplimiento} resume el cumplimiento de cada objetivo frente a su criterio
de aceptación y la evidencia que lo sustenta.

\begin{table}[H]
\centering
\small
\begin{tabular}{|L{2.4cm}|L{1.9cm}|L{8.5cm}|}
\hline
\textbf{Objetivo} & \textbf{Estado} & \textbf{Evidencia principal} \\ \hline
OE1 & Cumplido & Cuatro vectores mapeados a OWASP Top 10 2021 y MITRE ATT\&CK con fuente oficial; cinco clases canónicas que estructuran el dataset y la evaluación. \\ \hline
OE2 & Cumplido & 383{,}202 muestras etiquetadas (3.8 veces el mínimo) de siete fuentes auditadas, dos listas públicas y conjuntos propios, con etiquetado determinista y partición verificada. \\ \hline
OE3 & Cumplido, con reservas en la cobertura & F1 de 0.9239 a 0.9956 en 4 de 4 categorías; Isolation Forest con 91.25~\% de detección y 5.46~\% de falsos positivos; cobertura sobre payloads no vistos de 92.5 a 100~\% en las cuatro categorías con el paquete vigente 0.1.1 y el ejecutor corregido (Path Traversal 200/200, intervalo 98.1--100.0~\%; XSS 200/200, 98.1--100.0~\%). XSS: cumplido en cobertura (n=200); la explotabilidad confirmada sigue siendo n=1. Reservas: la medición de Path Traversal y de Command Injection se hizo sobre el corpus que reveló los defectos de medición, y la de XSS sobre un corpus construido después de conocer los resultados de la Ronda~5. \\ \hline
Objetivo general & Cumplido de forma parcial & Librería publicada en npm; variantes no vistas sobre el criterio de 80~\% con el Random Forest; extracción de 0.0273~ms y $\Delta p95$ de 12.08 y 12.85~ms. No se cumplen memoria (+332~MiB), CPU (+89~pp) ni la carga de volumen; el aporte del Isolation Forest a las variantes solo tiene evidencia indirecta (19 de 39 en validación). \\ \hline
\end{tabular}
\caption[Cumplimiento de los objetivos]{Cumplimiento de los objetivos frente a sus criterios de aceptación}
\label{tab:dis-cumplimiento}
\end{table}
{\footnotesize Nota. Elaboración propia con base en el Capítulo~\ref{cap:resultados}.}

% ------------------------------------------------------------------------------
\section{Hipótesis y validación de la solución}
\label{sec:dis-validacion}

El protocolo y el capítulo de Objetivos no formulan una hipótesis estadística contrastable, de
modo que no corresponde aceptarla o rechazarla. El objetivo general cumple ese papel como
afirmación de factibilidad: que un middleware de aprendizaje automático distribuido como paquete
npm puede detectar y clasificar ataques web sin comprometer el rendimiento del servidor anfitrión.
La evidencia respalda la primera parte (la detección y clasificación) y la mitad de la segunda
(la latencia en carga normal y de ataque), y no respalda el resto: los recursos de memoria y CPU,
el escenario de volumen y el comportamiento con tráfico legítimo real sin calibración previa. La
afirmación queda, por tanto, respaldada de forma parcial.

La validación de la solución combinó cuatro niveles de prueba. A nivel funcional, pruebas
automatizadas en integración continua (incluida la paridad numérica entre Python y ONNX) y la
instalación del paquete empaquetado en un directorio limpio contra una aplicación Express real.
A nivel de clasificación, la lectura del conjunto de prueba bloqueado por hash. A nivel de sistema,
las cuatro configuraciones de defensa sobre la aplicación de referencia, con un corpus filtrado
contra el entrenamiento, y el benchmark de rendimiento en Linux nativo con cinco repeticiones por
condición. Y a nivel externo, la integración en aplicaciones de terceros. Los criterios de
aceptación se aplicaron tal como los fijaron el Alcance y la Metodología, y el
Cuadro~\ref{tab:dis-cumplimiento} registra el estado de cada uno. Los dos niveles que más
información aportaron fueron el corpus
filtrado, que redujo la cobertura atribuible de Command Injection de 100~\% a 91.5~\%, y las
aplicaciones de terceros, una validación complementaria, que mostraron que las métricas del corpus no predicen el comportamiento
con tráfico legítimo real. Dos explicaciones quedan como posibles y no comprobadas: que el sesgo de
la cabecera \texttt{User-Agent} sea la causa principal de los bloqueos indebidos, y que el costo de
memoria y CPU se deba a la arquitectura de hilos y no a los modelos.

% ------------------------------------------------------------------------------
\section{Limitaciones}
\label{sec:dis-limitaciones}

Las amenazas a la validez descritas en la Sección~\ref{sec:met-validez} acotan el alcance de estas
conclusiones. El conjunto de prueba se volvió a bloquear en cada generación del modelo y se leyó
tres veces en la generación vigente; el corpus tiene una saturación de casi-duplicados que puede
inflar las métricas, sobre todo las de SQL Injection (aproximadamente 395 veces más pares casi
duplicados que filas, en una clase que es el 59.3~\% del conjunto de prueba) y que no se
cuantificó; el origen de los datos restaurados no está registrado, y el
17.8~\% de las filas de Command Injection del conjunto de prueba proviene de conjuntos sintéticos
sin generador versionado. Las evaluaciones usaron tres pares de modelos con la misma arquitectura
de ejecución pero no con la misma entrada (67 y 61 características en el par A; 69 y 63 en los otros dos), y el paquete con el que se hizo la medición principal de cobertura y se probaron las aplicaciones de terceros, el 0.1.0, no contiene el modelo cuyas
métricas de prueba se reportan, que solo está en la versión 0.1.1, posterior a esas mediciones.
La cobertura del modelo vigente sobre el corpus limpio no proviene de una evaluación estrictamente
ciega, y cada medición de cobertura corresponde a una sola corrida. La cobertura de
Path Traversal y de Command Injection con el paquete 0.1.1 se midió sobre el corpus que reveló los
defectos; las mediciones del paquete 0.1.0 y del tarball de validación se hicieron con el ejecutor
anterior y no se repitieron; la cobertura de XSS se midió sobre un corpus construido después de
conocer los resultados de la Ronda~5 y la explotabilidad confirmada descansa en un único caso; y
quedó pendiente la validación con una startup en producción. Estas
limitaciones no modifican las cifras reportadas, pero cada una indica una verificación que
fortalecería la evidencia en trabajos posteriores.

% ------------------------------------------------------------------------------
\section{Sugerencias derivadas de los resultados}
\label{sec:dis-sugerencias}

Los resultados sugieren las siguientes líneas, que se presentan como indicios que dan paso a las
conclusiones y a las recomendaciones, no como compromisos ni instrucciones. Cada una indica la
evidencia de la que parte.
\begin{itemize}
    \item \textbf{Instalar primero en modo \texttt{monitor}.} Los bloqueos de tráfico legítimo en
    aplicaciones de terceros (92.5~\% en dvws-node con un \texttt{User-Agent} de navegador) sugieren
    que la calibración con el tráfico propio debería preceder al bloqueo.
    \item \textbf{Diversificar el corpus legítimo.} El 0.4~\% de filas legítimas con un
    \texttt{User-Agent} largo, frente a hasta 97.7~\% en las clases de ataque, sugiere que la
    representación de clientes HTTP es la primera variable que conviene probar.
    \item \textbf{Medir el costo por componente.} Un incremento casi constante de 332~MiB entre
    escenarios sugiere un costo fijo de la arquitectura, pero no permite repartirlo entre el Random
    Forest, el Isolation Forest y el entorno de ejecución.
    \item \textbf{Proteger al servidor ante saturación.} El patrón de tiempos de espera al final de
    cada repetición del escenario de volumen sugiere una cola que crece; un mecanismo de
    \textit{fail-open} por longitud de cola podría evitarlo.
    \item \textbf{Repetir la medición con un corpus nuevo.} Que el 100~\% de Path Traversal se
    obtuviera sobre el corpus que reveló los defectos sugiere confirmarlo con otra semilla antes de
    atribuirlo a la librería.
    \item \textbf{Aislar la contribución del Isolation Forest.} Que marque 19 de 39 ataques que el
    Random Forest omite, a un costo de 866 bloqueos indebidos, sugiere evaluarlo fuera del conjunto
    de validación y, quizá, con un umbral distinto.
\end{itemize}

\ifdefined\modoensamble
\else
	\section*{Referencias}
	\printbibliography[heading=none]
	\end{document}
\fi
